\section{Proj of a graded ring}
\label{section-proj}

\noindent
In this section we construct Proj of a graded ring
following \cite[II, Section 2]{EGA}.

\medskip\noindent
Let $S$ be a graded ring. Consider the topological space $\text{Proj}(S)$
associated to $S$, see Algebra, Section \ref{algebra-section-proj}.
We will endow this space with a sheaf of rings $\mathcal{O}_{\text{Proj}(S)}$
such that the resulting pair $(\text{Proj}(S), \mathcal{O}_{\text{Proj}(S)})$
will be a scheme.

\medskip\noindent
Recall that $\text{Proj}(S)$ has a basis of open sets $D_{+}(f)$,
$f \in S_d$, $d \geq 1$ which we call {\it standard opens}, see Algebra,
Section \ref{algebra-section-proj}. This terminology will always
imply that $f$ is homogeneous of positive degree even if we forget to
mention it. In addition, the intersection of two standard opens is another:
$D_{+}(f) \cap D_{+}(g) = D_{+}(fg)$, for $f, g \in S$ homogeneous of positive
degree.

\begin{lemma}
\label{lemma-standard-open}
Let $S$ be a graded ring. Let $f \in S$ homogeneous of positive degree.
\begin{enumerate}
\item If $g\in S$ homogeneous of positive degree
and $D_{+}(g) \subset D_{+}(f)$, then
\begin{enumerate}
\item $f$ is invertible in $S_g$, and
$f^{\deg(g)}/g^{\deg(f)}$ is invertible in $S_{(g)}$,
\item $g^e = af$ for some $e \geq 1$ and $a \in S$ homogeneous,
\item there is a canonical $S$-algebra map $S_f \to S_g$,
\item there is a canonical $S_0$-algebra map $S_{(f)} \to S_{(g)}$
compatible with the map $S_f \to S_g$,
\item the map $S_{(f)} \to S_{(g)}$ induces an isomorphism
$$
(S_{(f)})_{g^{\deg(f)}/f^{\deg(g)}} \cong S_{(g)},
$$
\item these maps induce a commutative diagram of
topological spaces
$$
\xymatrix{
D_{+}(g) \ar[d] &
\{\mathbf{Z}\text{-graded primes of }S_g\} \ar[l] \ar[r] \ar[d] &
\Spec(S_{(g)}) \ar[d] \\
D_{+}(f) &
\{\mathbf{Z}\text{-graded primes of }S_f\} \ar[l] \ar[r] &
\Spec(S_{(f)})
}
$$
where the horizontal maps are homeomorphisms and the vertical maps
are open immersions,
\item there are compatible canonical $S_f$ and $S_{(f)}$-module
maps $M_f \to M_g$ and $M_{(f)} \to M_{(g)}$ for any graded $S$-module $M$,
and
\item the map $M_{(f)} \to M_{(g)}$ induces an isomorphism
$$
(M_{(f)})_{g^{\deg(f)}/f^{\deg(g)}} \cong M_{(g)}.
$$
\end{enumerate}
\item Any open covering of $D_{+}(f)$ can be refined to a finite
open covering of the form $D_{+}(f) = \bigcup_{i = 1}^n D_{+}(g_i)$.
\item Let $g_1, \ldots, g_n \in S$ be homogeneous of positive degree.
Then $D_{+}(f) \subset \bigcup D_{+}(g_i)$
if and only if
$g_1^{\deg(f)}/f^{\deg(g_1)}, \ldots, g_n^{\deg(f)}/f^{\deg(g_n)}$
generate the unit ideal in $S_{(f)}$.
\end{enumerate}
\end{lemma}

\begin{proof}
Recall that $D_{+}(g) = \Spec(S_{(g)})$ with identification
given by the ring maps $S \to S_g \leftarrow S_{(g)}$, see
Algebra, Lemma \ref{algebra-lemma-topology-proj}.
Thus $f^{\deg(g)}/g^{\deg(f)}$ is an element of $S_{(g)}$ which is not
contained in any prime ideal, and hence invertible,
see Algebra, Lemma \ref{algebra-lemma-Zariski-topology}.
We conclude that (a) holds.
Write the inverse of $f$ in $S_g$ as $a/g^d$.
We may replace $a$ by its homogeneous part of degree $d\deg(g) - \deg(f)$.
This means $g^d - af$ is annihilated by a power of $g$, whence
$g^e = af$ for some $a \in S$ homogeneous of degree $e\deg(g) - \deg(f)$.
This proves (b).
For (c), the map $S_f \to S_g$ exists by (a) from the universal property
of localization, or we can define it by mapping $b/f^n$
to $a^nb/g^{ne}$. This clearly induces a map of the subrings
$S_{(f)} \to S_{(g)}$ of degree zero elements as well.
We can similarly define $M_f \to M_g$ and $M_{(f)} \to M_{(g)}$ by mapping
$x/f^n$ to $a^nx/g^{ne}$. The statements writing $S_{(g)}$
resp.\ $M_{(g)}$ as principal localizations of $S_{(f)}$ resp.\ $M_{(f)}$
are clear from the formulas above. The maps in the commutative diagram
of topological spaces correspond to the ring maps given above. The
horizontal arrows are homeomorphisms by
Algebra, Lemma \ref{algebra-lemma-topology-proj}.
The vertical arrows are open immersions since the left
one is the inclusion of an open subset.

\medskip\noindent
The open $D_{+}(f)$ is quasi-compact because it is homeomorphic
to $\Spec(S_{(f)})$, see  Algebra, Lemma \ref{algebra-lemma-quasi-compact}.
Hence the second statement follows directly
from the fact that the standard opens form
a basis for the topology.

\medskip\noindent
The third statement follows directly from
Algebra, Lemma \ref{algebra-lemma-Zariski-topology}.
\end{proof}

\noindent
In Sheaves, Section \ref{sheaves-section-bases} we defined
the notion of a sheaf on a basis, and we showed that it is
essentially equivalent to the notion of a sheaf on the space,
see Sheaves, Lemmas \ref{sheaves-lemma-extend-off-basis} and
\ref{sheaves-lemma-extend-off-basis-structures}. Moreover,
we showed in
Sheaves, Lemma \ref{sheaves-lemma-cofinal-systems-coverings-standard-case}
that it is sufficient to check the sheaf
condition on a cofinal system of open coverings for each
standard open. By the lemma above it suffices to check
on the finite coverings by standard opens.

\begin{definition}
\label{definition-standard-covering}
Let $S$ be a graded ring.
Suppose that $D_{+}(f) \subset \text{Proj}(S)$ is a standard
open. A {\it standard open covering} of $D_{+}(f)$
is a covering $D_{+}(f) = \bigcup_{i = 1}^n D_{+}(g_i)$,
where $g_1, \ldots, g_n \in S$ are homogeneous of positive degree.
\end{definition}

\noindent
Let $S$ be a graded ring. Let $M$ be a graded $S$-module. We will define
a presheaf $\widetilde M$ on the basis of standard opens.
Suppose that $U \subset \text{Proj}(S)$ is a standard open.
If $f, g \in S$ are homogeneous of positive degree
such that $D_{+}(f) = D_{+}(g)$, then
by Lemma \ref{lemma-standard-open} above there are canonical
maps $M_{(f)} \to M_{(g)}$ and $M_{(g)} \to M_{(f)}$ which are
mutually inverse. Hence we may choose any $f$ such that $U = D_{+}(f)$
and define
$$
\widetilde M(U) = M_{(f)}.
$$
Note that if $D_{+}(g) \subset D_{+}(f)$, then by
Lemma \ref{lemma-standard-open} above we have
a canonical map
$$
\widetilde M(D_{+}(f)) = M_{(f)} \longrightarrow
M_{(g)} = \widetilde M(D_{+}(g)).
$$
Clearly, this defines a presheaf of abelian groups on the basis
of standard opens. If $M = S$, then $\widetilde S$ is a presheaf
of rings on the basis of standard opens. And for general $M$ we
see that $\widetilde M$ is a presheaf of $\widetilde S$-modules
on the basis of standard opens.

\medskip\noindent
Let us compute the stalk of $\widetilde M$ at a point
$x \in \text{Proj}(S)$.
Suppose that $x$ corresponds to the homogeneous prime
ideal $\mathfrak p \subset S$.
By definition of the stalk we see that
$$
\widetilde M_x
=
\colim_{f\in S_d, d > 0, f\not\in \mathfrak p} M_{(f)}
$$
Here the set $\{f \in S_d, d > 0, f \not \in \mathfrak p\}$ is preordered by
the rule $f \geq f' \Leftrightarrow D_{+}(f) \subset D_{+}(f')$.
If $f_1, f_2 \in S \setminus \mathfrak p$ are homogeneous of positive
degree, then we have
$f_1f_2 \geq f_1$ in this ordering. In
Algebra, Section \ref{algebra-section-proj}
we defined $M_{(\mathfrak p)}$ as the module whose elements are fractions
$x/f$ with $x, f$ homogeneous, $\deg(x) = \deg(f)$, $f \not \in \mathfrak p$.
Since $\mathfrak p \in \text{Proj}(S)$ there exists at least
one $f_0 \in S$ homogeneous of positive degree with $f_0 \not\in \mathfrak p$.
Hence $x/f = f_0x/ff_0$ and we see that we may always assume
the denominator of an element in $M_{(\mathfrak p)}$ has positive degree.
From these remarks it follows easily that
$$
\widetilde M_x = M_{(\mathfrak p)}.
$$

\medskip\noindent
Next, we check the sheaf condition for the standard open coverings.
If $D_{+}(f) = \bigcup_{i = 1}^n D_{+}(g_i)$, then the sheaf condition
for this covering is equivalent with the exactness of the
sequence
$$
0 \to M_{(f)} \to \bigoplus M_{(g_i)} \to \bigoplus M_{(g_ig_j)}.
$$
Note that $D_{+}(g_i) = D_{+}(fg_i)$, and hence we can rewrite this
sequence as the sequence
$$
0 \to M_{(f)} \to \bigoplus M_{(fg_i)} \to \bigoplus M_{(fg_ig_j)}.
$$
By Lemma \ref{lemma-standard-open} we see that
$g_1^{\deg(f)}/f^{\deg(g_1)}, \ldots, g_n^{\deg(f)}/f^{\deg(g_n)}$
generate the unit ideal in $S_{(f)}$, and that the modules
$M_{(fg_i)}$, $M_{(fg_ig_j)}$ are the principal localizations
of the $S_{(f)}$-module $M_{(f)}$ at these elements and their products.
Thus we may apply Algebra, Lemma \ref{algebra-lemma-cover-module}
to the module $M_{(f)}$ over $S_{(f)}$ and the elements
$g_1^{\deg(f)}/f^{\deg(g_1)}, \ldots, g_n^{\deg(f)}/f^{\deg(g_n)}$.
We conclude that the sequence is exact. By the remarks
made above, we see that $\widetilde M$ is a sheaf
on the basis of standard opens.

\medskip\noindent
Thus we conclude from the material in
Sheaves, Section \ref{sheaves-section-bases}
that there exists a
unique sheaf of rings $\mathcal{O}_{\text{Proj}(S)}$
which agrees with $\widetilde S$ on the standard opens.
Note that by our computation of stalks above and
Algebra, Lemma \ref{algebra-lemma-proj-prime} the
stalks of this sheaf of rings are all local rings.

\medskip\noindent
Similarly, for any graded $S$-module $M$ there exists
a unique sheaf of $\mathcal{O}_{\text{Proj}(S)}$-modules
$\mathcal{F}$ which agrees with $\widetilde M$ on the
standard opens, see
Sheaves, Lemma \ref{sheaves-lemma-extend-off-basis-module}.

\begin{definition}
\label{definition-structure-sheaf}
Let $S$ be a graded ring.
\begin{enumerate}
\item The {\it structure sheaf $\mathcal{O}_{\text{Proj}(S)}$ of the
homogeneous spectrum of $S$} is the unique sheaf of rings
$\mathcal{O}_{\text{Proj}(S)}$
which agrees with $\widetilde S$ on the basis of standard opens.
\item The locally ringed space
$(\text{Proj}(S), \mathcal{O}_{\text{Proj}(S)})$ is called
the {\it homogeneous spectrum} of $S$ and denoted $\text{Proj}(S)$.
\item The sheaf of $\mathcal{O}_{\text{Proj}(S)}$-modules
extending $\widetilde M$ to all opens of $\text{Proj}(S)$
is called the sheaf of $\mathcal{O}_{\text{Proj}(S)}$-modules
associated to $M$. This sheaf is denoted $\widetilde M$ as
well.
\end{enumerate}
\end{definition}

\noindent
We summarize the results obtained so far.

\begin{lemma}
\label{lemma-proj-sheaves}
Let $S$ be a graded ring. Let $M$ be a graded $S$-module.
Let $\widetilde M$ be the sheaf of $\mathcal{O}_{\text{Proj}(S)}$-modules
associated to $M$.
\begin{enumerate}
\item For every $f \in S$ homogeneous of positive degree we have
$$
\Gamma(D_{+}(f), \mathcal{O}_{\text{Proj}(S)}) = S_{(f)}.
$$
\item For every $f\in S$ homogeneous of positive degree
we have $\Gamma(D_{+}(f), \widetilde M) = M_{(f)}$
as an $S_{(f)}$-module.
\item Whenever $D_{+}(g) \subset D_{+}(f)$ the restriction mappings
on $\mathcal{O}_{\text{Proj}(S)}$ and $\widetilde M$
are the maps
$S_{(f)} \to S_{(g)}$ and $M_{(f)} \to M_{(g)}$ from Lemma
\ref{lemma-standard-open}.
\item Let $\mathfrak p$ be a homogeneous prime of $S$ not containing
$S_{+}$, and let $x \in \text{Proj}(S)$
be the corresponding point. We have
$\mathcal{O}_{\text{Proj}(S), x} = S_{(\mathfrak p)}$.
\item Let $\mathfrak p$ be a homogeneous prime of $S$ not containing
$S_{+}$, and let $x \in \text{Proj}(S)$
be the corresponding point. We have $(\widetilde M)_x = M_{(\mathfrak p)}$
as an $S_{(\mathfrak p)}$-module.
\item
\label{item-map}
There is a canonical ring map
$
S_0 \longrightarrow \Gamma(\text{Proj}(S), \widetilde S)
$
and a canonical $S_0$-module map
$
M_0 \longrightarrow \Gamma(\text{Proj}(S), \widetilde M)
$
compatible with the descriptions of sections over standard opens
and stalks above.
\end{enumerate}
Moreover, all these identifications are functorial in the graded
$S$-module $M$. In particular, the functor $M \mapsto \widetilde M$
is an exact functor from the category of graded $S$-modules
to the category of $\mathcal{O}_{\text{Proj}(S)}$-modules.
\end{lemma}

\begin{proof}
Assertions (1) - (5) are clear from the discussion above.
We see (6) since there are canonical maps $M_0 \to M_{(f)}$,
$x \mapsto x/1$ compatible with the restriction maps
described in (3). The exactness of the functor $M \mapsto \widetilde M$
follows from the fact that the functor $M \mapsto M_{(\mathfrak p)}$
is exact (see Algebra, Lemma \ref{algebra-lemma-proj-prime})
and the fact that exactness of short exact sequences
may be checked on stalks, see
Modules, Lemma \ref{modules-lemma-abelian}.
\end{proof}

\begin{remark}
\label{remark-global-sections-not-isomorphism}
The map from $M_0$ to the global sections of $\widetilde M$
is generally far from being an isomorphism. A trivial
example is to take $S = k[x, y, z]$ with $1 = \deg(x) = \deg(y) = \deg(z)$
(or any number of variables) and to take $M = S/(x^{100}, y^{100}, z^{100})$.
It is easy to see that $\widetilde M = 0$, but $M_0 = k$.
\end{remark}

\begin{lemma}
\label{lemma-standard-open-proj}
Let $S$ be a graded ring. Let $f \in S$ be homogeneous of positive degree.
Suppose that $D(g) \subset \Spec(S_{(f)})$ is a standard open.
Then there exists an $h \in S$ homogeneous of positive degree such that
$D(g)$ corresponds to $D_{+}(h) \subset D_{+}(f)$ via the homeomorphism
of Algebra, Lemma \ref{algebra-lemma-topology-proj}. In fact we can
take $h$ such that $g = h/f^n$ for some $n$.
\end{lemma}

\begin{proof}
Write $g = h/f^n$ for some $h$ homogeneous of positive degree
and some $n \geq 1$. If $D_{+}(h)$ is not contained in
$D_{+}(f)$ then we replace $h$ by $hf$ and $n$ by $n + 1$.
Then $h$ has the required shape and $D_{+}(h) \subset D_{+}(f)$
corresponds to $D(g) \subset \Spec(S_{(f)})$.
\end{proof}

\begin{lemma}
\label{lemma-proj-scheme}
Let $S$ be a graded ring.
The locally ringed space $\text{Proj}(S)$ is a scheme.
The standard opens $D_{+}(f)$ are affine opens.
For any graded $S$-module $M$ the sheaf
$\widetilde M$ is a quasi-coherent sheaf of
$\mathcal{O}_{\text{Proj}(S)}$-modules.
\end{lemma}

\begin{proof}
Consider a standard open $D_{+}(f) \subset \text{Proj}(S)$.
By Lemmas \ref{lemma-standard-open} and \ref{lemma-proj-sheaves}
we have $\Gamma(D_{+}(f), \mathcal{O}_{\text{Proj}(S)}) = S_{(f)}$, and
we have a homeomorphism $\varphi : D_{+}(f) \to \Spec(S_{(f)})$.
For any standard open $D(g) \subset \Spec(S_{(f)})$ we may
pick an $h \in S_{+}$ as in Lemma \ref{lemma-standard-open-proj}.
Then $\varphi^{-1}(D(g)) = D_{+}(h)$, and by
Lemmas \ref{lemma-proj-sheaves} and \ref{lemma-standard-open} we see
$$
\Gamma(D_{+}(h), \mathcal{O}_{\text{Proj}(S)})
=
S_{(h)}
=
(S_{(f)})_{h^{\deg(f)}/f^{\deg(h)}}
=
(S_{(f)})_g
=
\Gamma(D(g), \mathcal{O}_{\Spec(S_{(f)})}).
$$
Thus the restriction of $\mathcal{O}_{\text{Proj}(S)}$ to
$D_{+}(f)$ corresponds via the homeomorphism $\varphi$
exactly to the sheaf $\mathcal{O}_{\Spec(S_{(f)})}$
as defined in Schemes, Section \ref{schemes-section-affine-schemes}.
We conclude that $D_{+}(f)$ is an affine scheme isomorphic to
$\Spec(S_{(f)})$ via $\varphi$ and
hence that $\text{Proj}(S)$ is a scheme.

\medskip\noindent
In exactly the same way we show that $\widetilde M$ is a
quasi-coherent sheaf of $\mathcal{O}_{\text{Proj}(S)}$-modules.
Namely, the argument above will show that
$$
\widetilde M|_{D_{+}(f)} \cong \varphi^*\left(\widetilde{M_{(f)}}\right)
$$
which shows that $\widetilde M$ is quasi-coherent.
\end{proof}

\begin{lemma}
\label{lemma-proj-separated}
Let $S$ be a graded ring.
The scheme $\text{Proj}(S)$ is separated.
\end{lemma}

\begin{proof}
We have to show that the canonical morphism
$\text{Proj}(S) \to \Spec(\mathbf{Z})$
is separated.
We will use Schemes, Lemma \ref{schemes-lemma-characterize-separated}.
Thus it suffices to show given any pair of standard opens
$D_{+}(f)$ and $D_{+}(g)$ that $D_{+}(f) \cap D_{+}(g) = D_{+}(fg)$
is affine (clear) and that the ring map
$$
S_{(f)} \otimes_{\mathbf{Z}} S_{(g)} \longrightarrow S_{(fg)}
$$
is surjective. Any element $s$ in $S_{(fg)}$ is of
the form $s = h/(f^ng^m)$ with $h \in S$ homogeneous of degree
$n\deg(f) + m\deg(g)$. We may multiply $h$ by a suitable
monomial $f^ig^j$ and assume that $n = n' \deg(g)$, and
$m = m' \deg(f)$. Then we can rewrite $s$ as
$s = h/f^{(n' + m')\deg(g)} \cdot f^{m'\deg(g)}/g^{m'\deg(f)}$.
So $s$ is indeed in the image of the displayed arrow.
\end{proof}

\begin{lemma}
\label{lemma-proj-quasi-compact}
Let $S$ be a graded ring.
The scheme $\text{Proj}(S)$ is quasi-compact if and only
if there exist finitely many homogeneous elements
$f_1, \ldots, f_n \in S_{+}$ such that
$S_{+} \subset \sqrt{(f_1, \ldots, f_n)}$. In this case
$\text{Proj}(S) = D_+(f_1) \cup \ldots \cup D_+(f_n)$.
\end{lemma}

\begin{proof}
Given such a collection of elements the standard affine opens
$D_{+}(f_i)$ cover $\text{Proj}(S)$ by
Algebra, Lemma \ref{algebra-lemma-topology-proj}.
Conversely, if $\text{Proj}(S)$ is quasi-compact, then we
may cover it by finitely many standard opens
$D_{+}(f_i)$, $i = 1, \ldots, n$ and we see that
$S_{+} \subset \sqrt{(f_1, \ldots, f_n)}$ by the
lemma referenced above.
\end{proof}

\begin{lemma}
\label{lemma-structure-morphism-proj}
Let $S$ be a graded ring. The scheme $\text{Proj}(S)$ has a canonical morphism
towards the affine scheme $\Spec(S_0)$, agreeing with the map on
topological spaces coming from
Algebra, Definition \ref{algebra-definition-proj}.
\end{lemma}

\begin{proof}
We saw above that our construction of $\widetilde S$,
resp.\ $\widetilde M$ gives a sheaf of $S_0$-algebras, resp.\ $S_0$-modules.
Hence we get a morphism by
Schemes, Lemma \ref{schemes-lemma-morphism-into-affine}.
This morphism, when restricted to $D_{+}(f)$ comes from the
canonical ring map $S_0 \to S_{(f)}$. The maps
$S \to S_f$, $S_{(f)} \to S_f$ are $S_0$-algebra maps, see
Lemma \ref{lemma-standard-open}.
Hence if the homogeneous prime $\mathfrak p \subset S$
corresponds to the $\mathbf{Z}$-graded prime $\mathfrak p' \subset S_f$
and the (usual) prime $\mathfrak p'' \subset S_{(f)}$, then
each of these has the same inverse image in $S_0$.
\end{proof}

\begin{lemma}
\label{lemma-proj-valuative-criterion}
Let $S$ be a graded ring. If $S$ is finitely generated as
an algebra over $S_0$, then
the morphism $\text{Proj}(S) \to \Spec(S_0)$ satisfies
the existence and uniqueness parts of the valuative criterion,
see Schemes, Definition \ref{schemes-definition-valuative-criterion}.
\end{lemma}

\begin{proof}
The uniqueness part follows from the fact that $\text{Proj}(S)$ is
separated (Lemma \ref{lemma-proj-separated} and
Schemes, Lemma \ref{schemes-lemma-separated-implies-valuative}).
Choose $x_i \in S_{+}$ homogeneous, $i = 1, \ldots, n$
which generate $S$ over $S_0$. Let $d_i = \deg(x_i)$ and
set $d = \text{lcm}\{d_i\}$. Suppose we are given a diagram
$$
\xymatrix{
\Spec(K) \ar[r] \ar[d] & \text{Proj}(S) \ar[d] \\
\Spec(A) \ar[r] & \Spec(S_0)
}
$$
as in Schemes, Definition \ref{schemes-definition-valuative-criterion}.
Denote $v : K^* \to \Gamma$ the valuation of $A$, see
Algebra, Definition \ref{algebra-definition-value-group}.
We may choose an $f \in S_{+}$ homogeneous such that
$\Spec(K)$ maps into $D_{+}(f)$. Then we get a commutative
diagram of ring maps
$$
\xymatrix{
K & S_{(f)} \ar[l]^{\varphi} \\
A \ar[u] & S_0 \ar[l] \ar[u]
}
$$
After renumbering we may assume that $\varphi(x_i^{\deg(f)}/f^{d_i})$
is nonzero for $i = 1, \ldots, r$ and zero for $i = r + 1, \ldots, n$.
Since the open sets $D_{+}(x_i)$ cover $\text{Proj}(S)$ we see that $r \geq 1$.
Let $i_0 \in \{1, \ldots, r\}$ be an index minimizing
$\gamma_i = (d/d_i)v(\varphi(x_i^{\deg(f)}/f^{d_i}))$ in $\Gamma$.
For convenience set $x_0 = x_{i_0}$ and $d_0 = d_{i_0}$.
The ring map $\varphi$ factors though a map $\varphi' : S_{(fx_0)} \to K$
which gives a ring map $S_{(x_0)} \to S_{(fx_0)} \to K$.
The algebra $S_{(x_0)}$ is generated over $S_0$ by the elements
$x_1^{e_1} \ldots x_n^{e_n}/x_0^{e_0}$, where $\sum e_i d_i = e_0 d_0$.
If $e_i > 0$ for some $i > r$, then
$\varphi'(x_1^{e_1} \ldots x_n^{e_n}/x_0^{e_0}) = 0$.
If $e_i = 0$ for $i > r$, then we have
\begin{align*}
d \deg(f) v(\varphi'(x_1^{e_1} \ldots x_r^{e_r}/x_0^{e_0}))
& =
d v(\varphi'(x_1^{e_1 \deg(f)} \ldots x_r^{e_r \deg(f)}/x_0^{e_0 \deg(f)})) \\
& =
d \sum e_i v(\varphi'(x_i^{\deg(f)}/f^{d_i}))
- e_0 v(\varphi'(x_0^{\deg(f)}/f^{d_0})) \\
& =
\sum e_i d_i \gamma_i - e_0 d_0 \gamma_0 \\
& \geq
\sum e_i d_i \gamma_0 - e_0 d_0 \gamma_0 = 0
\end{align*}
because $\gamma_0$ is minimal among the $\gamma_i$.
This implies that $S_{(x_0)}$ maps into $A$ via $\varphi'$.
The corresponding morphism of schemes
$\Spec(A) \to \Spec(S_{(x_0)}) = D_{+}(x_0)
\subset \text{Proj}(S)$ provides the morphism fitting into
the first commutative diagram of this proof.
\end{proof}

\noindent
We saw in the proof of Lemma \ref{lemma-proj-valuative-criterion}
that, under the hypotheses of that lemma, the morphism
$\text{Proj}(S) \to \Spec(S_0)$ is quasi-compact as well. Hence (by
Schemes, Proposition \ref{schemes-proposition-characterize-universally-closed})
we see that $\text{Proj}(S) \to \Spec(S_0)$ is universally closed in
the situation of the lemma. We give several examples showing these results
do not hold without some assumption on the graded ring $S$.

\begin{example}
\label{example-not-existence-valuative-big-proj}
Let $\mathbf{C}[X_1, X_2, X_3, \ldots]$ be the graded $\mathbf{C}$-algebra
with each $X_i$ in degree $1$. Consider the ring map
$$
\mathbf{C}[X_1, X_2, X_3, \ldots]
\longrightarrow
\mathbf{C}[t^\alpha ; \alpha \in \mathbf{Q}_{\geq 0}]
$$
which maps $X_i$ to $t^{1/i}$. The right hand side becomes a valuation ring
$A$ upon localization at the ideal $\mathfrak m = (t^\alpha ; \alpha > 0)$.
Let $K$ be the fraction field of $A$. The above gives a morphism
$\Spec(K) \to \text{Proj}(\mathbf{C}[X_1, X_2, X_3, \ldots])$ which does not
extend to a morphism defined on all of $\Spec(A)$.
The reason is that the image of $\Spec(A)$ would be contained
in one of the $D_{+}(X_i)$ but then $X_{i + 1}/X_i$ would map
to an element of $A$ which it doesn't since it maps to
$t^{1/(i + 1) - 1/i}$.
\end{example}

\begin{example}
\label{example-not-existence-valuative-small-proj}
Let $R = \mathbf{C}[t]$ and
$$
S = R[X_1, X_2, X_3, \ldots]/(X_i^2 - tX_{i + 1}).
$$
The grading is such that $R = S_0$ and $\deg(X_i) = 2^{i - 1}$.
Note that if $\mathfrak p \in \text{Proj}(S)$ then
$t \not \in \mathfrak p$ (otherwise $\mathfrak p$ has to contain
all of the $X_i$ which is not allowed for an element of
the homogeneous spectrum). Thus we see that
$D_{+}(X_i) = D_{+}(X_{i + 1})$ for all $i$. Hence
$\text{Proj}(S)$ is quasi-compact; in fact it is affine
since it is equal to $D_{+}(X_1)$. It is easy to see that
the image of $\text{Proj}(S) \to \Spec(R)$ is
$D(t)$. Hence the morphism $\text{Proj}(S) \to \Spec(R)$
is not closed. Thus the valuative criterion cannot apply because
it would imply that the morphism is closed (see
Schemes, Proposition \ref{schemes-proposition-characterize-universally-closed}
).
\end{example}

\begin{example}
\label{example-trivial-proj}
Let $A$ be a ring.
Let $S = A[T]$ as a graded $A$ algebra with $T$ in degree $1$.
Then the canonical morphism $\text{Proj}(S) \to \Spec(A)$
(see Lemma \ref{lemma-structure-morphism-proj})
is an isomorphism.
\end{example}

\begin{example}
\label{example-open-subset-proj}
Let $X = \Spec(A)$ be an affine scheme, and let $U \subset X$
be an open subscheme.  Grade $A[T]$ by setting $\deg T = 1$.  Define $S$
to be the subring of $A[T]$ generated by $A$ and all $fT^i$, where $i \ge 0$
and where $f \in A$ is such that $D(f) \subset U$.  We claim that $S$
is a graded ring with $S_0 = A$ such that $\text{Proj}(S) \cong U$,
and this isomorphism identifies the canonical morphism
$\text{Proj}(S) \to \Spec(A)$ of Lemma \ref{lemma-structure-morphism-proj}
with the inclusion $U \subset X$.

\medskip\noindent
Suppose $\mathfrak p \in \text{Proj}(S)$ is such that every $fT \in S_1$
is in $\mathfrak p$.  Then every generator $fT^i$ with $i \ge 1$
is in $\mathfrak p$ because $(fT^i)^2 = (fT)(fT^{2i-1}) \in \mathfrak p$
and $\mathfrak p$ is radical.  But then $\mathfrak p \supset S_+$, which
is impossible. Consequently $\text{Proj}(S)$ is covered by the standard
open affine subsets $\{D_+(fT)\}_{fT \in S_1}$.

\medskip\noindent
Observe that, if $fT \in S_1$, then the inclusion $S \subset A[T]$
induces a graded isomorphism of $S[(fT)^{-1}]$ with $A[T, T^{-1}, f^{-1}]$.
Hence the standard open subset $D_+(fT) \cong \Spec(S_{(fT)})$
is isomorphic to $\Spec(A[T, T^{-1}, f^{-1}]_0) = \Spec(A[f^{-1}])$.
It is clear that this isomorphism is a restriction of the canonical morphism
$\text{Proj}(S) \to \Spec(A)$.  If in addition $gT \in S_1$, then
$S[(fT)^{-1}, (gT)^{-1}] \cong A[T, T^{-1}, f^{-1}, g^{-1}]$
as graded rings, so $D_+(fT) \cap D_+(gT) \cong \Spec(A[f^{-1}, g^{-1}])$.
Therefore $\text{Proj}(S)$ is the union of open subschemes $D_+(fT)$
which are isomorphic to the open subschemes $D(f) \subset X$
under the canonical morphism, and these open subschemes intersect
in $\text{Proj}(S)$ in the same way they do in $X$.
We conclude that the canonical morphism is an isomorphism of
$\text{Proj}(S)$ with the union of all $D(f) \subset U$, which is $U$.
\end{example}





















